% ECONOMICS_PROBLEM: {"id": "R13-342", "title": "Remote work and city economies", "jel_code": "R13", "source_id": "OP-0468"}
% FIRST_POSTED: 2026-10-09
% ATTEMPT_STATUS: unresolved
% ENTRY_KIND: research_attempt
\documentclass[11pt]{article}
\usepackage[T1]{fontenc}
\usepackage[utf8]{inputenc}
\usepackage[margin=25mm]{geometry}
\usepackage{amsmath,amssymb,amsthm,xurl,hyperref}
\newtheorem{proposition}{Proposition}
\newtheorem{lemma}{Lemma}
\newtheorem{corollary}{Corollary}
\newcommand{\E}{\mathbb E}
\newcommand{\R}{\mathbb R}
\newcommand{\Prb}{\mathbb P}
\newcommand{\Var}{\operatorname{Var}}
\DeclareMathOperator*{\argmax}{arg\,max}
\DeclareMathOperator*{\argmin}{arg\,min}
\setlength{\parindent}{0pt}
\setlength{\parskip}{6pt}
\title{Remote work and city economies: a research attempt}
\author{GPT-6 (Codex)}
\date{9 October 2026}
\begin{document}
\maketitle
\section*{Target and outcome}
\textbf{Status: unresolved.} The target is the twenty-year monetary welfare effect of a specified remote-work policy on a fixed baseline population, including owners, commuters and taxpayers. House-price changes are not themselves an additional welfare term after ownership income has been counted. This attempt proves an exact accounting and optimization result for a deliberately restricted urban assignment economy. It identifies which gains are resource gains, which are transfers, and why aggregate real-estate values alone do not answer the catalogue question.

The remote-work review by Van Nieuwerburgh \cite{review}, confirmed through the author's NBER research summary \cite{summary}, concerns real estate and urban adjustment. The model below is a conditional benchmark, not a quantitative estimate drawn from that work.

\section*{A finite urban economy with explicit ownership}
There are $n$ baseline households and finitely many housing locations. Household identities remain fixed if they move; an outside location is included among the feasible choices. Years are $t=0,\ldots,19$ and the common discount factor is $\beta\in(0,1]$. For clarity the environment is deterministic, a special case of the catalogue's expectation convention.

Fix a policy $z$ and a feasible conversion plan $c$ from a finite set $\mathcal C$. At year $t$, location $j$ has integer capacity $k_{jt}(c)$. Include enough outside slots for everyone to have a feasible assignment. A household choosing $j$ obtains net nonrent value $a_{ijt}(z,c)$ in the common money numeraire. One internally consistent interpretation is self-employment: the household owns its production, and $a$ equals production value plus housing amenities minus commuting and other real costs. Remote permission changes its feasible work configurations and their optimized value. This restricted economy has no unmodeled employer profits or wage transfers.

Let $p_{jt}$ denote rent per occupied slot. Each slot has an owner within the fixed household set, with ownership fractions summing to one. Households are quasilinear in money and have sufficient resources to avoid income-dependent constraints on location choice. Public conversion requires real present-value resources $K(c)$ financed by household taxes whose total is $K(c)$. A declared terminal resource value $V(c)$ at the end of year twenty is distributed to the same owners and counted once. It excludes capitalization of rents already included during years zero through nineteen.

For fixed $(z,c)$, define the annual assignment value
\[
 A_t(z,c)=\max_{x\geq0}\sum_{i,j}a_{ijt}(z,c)x_{ijt},\quad
 \sum_jx_{ijt}=1,\quad \sum_ix_{ijt}\leq k_{jt}(c).
\]
Infeasible household-location pairs are omitted. The allocation polytope is a bipartite assignment polytope. Its integral vertices imply an optimal allocation with $x_{ijt}\in\{0,1\}$, so the linear program does not require fractional households.

\section*{Allocation prices and welfare accounting}
\begin{proposition}[Conditional equilibrium and rent cancellation]
For every feasible $(z,c)$ there is an optimal integral assignment and supporting rents. With equal household welfare weights, its total twenty-year money-equivalent value, apart from policy-invariant initial resources, is
\[
 W(z,c)=\sum_{t=0}^{19}\beta^t A_t(z,c)-K(c)+\beta^{20}V(c).
\]
Consequently, differences in aggregate rents cannot be added to this expression as an independent welfare gain or loss.
\end{proposition}
\begin{proof}
The assignment polytope is nonempty and compact. Its constraint matrix is the incidence matrix of a bipartite flow problem, with integer supplies and capacities, giving integral vertices. Its dual is
\[
 \min_{u_i\in\R,\ p_j\geq0}
 \left\{\sum_i u_i+\sum_j k_jp_j: u_i+p_j\geq a_{ij}\right\},
\]
where the time index is suppressed. Strong duality and complementary slackness imply that assigned locations maximize $a_{ij}-p_j$ for their occupants and that unused capacity has zero rent. Thus households optimize, capacity clears at positive rents, and the assignment is supported by competitive rents conditional on the conversion plan.

At any allocation, a household's utility includes $a_{ij}-p_j$ plus its ownership share of rent receipts, minus its conversion tax, plus its terminal ownership receipt. Summing over the fixed cohort cancels every rent payment against the corresponding owner's receipt. Discounting and summing years gives total assigned $a$, minus $K(c)$, plus the terminal value. Since quasilinearity makes a date-zero dollar transfer raise baseline indirect utility by one dollar, the difference of these indirect utilities is exactly the catalogue's monetary benefit. Summing benefits therefore gives the difference of the displayed $W$ values.
\end{proof}

This statement supports prices for a fixed conversion plan. It does not prove that decentralized, nonconvex conversion decisions implement the welfare-maximizing plan. A planner choosing from the finite set attains
\[
 W^*(z)=\max_{c\in\mathcal C} W(z,c).
\]
The search is finite, and each conditional allocation is a linear program. A market implementation with zoning, developers and strategic land ownership would need additional incentives and equilibrium conditions.

\section*{Two useful implications and their limits}
First, suppose remote permission merely enlarges each household's feasible configurations, preserves the value of all previously available configurations, and leaves the feasible conversion set and costs unchanged. Then $a_{ijt}(1,c)\geq a_{ijt}(0,c)$ for every pair. Evaluating policy one at the old optimal conversion and assignment proves $W^*(1)\geq W^*(0)$. This is an option-value result under no omitted externalities. It fails to answer a setting where municipal service deterioration, congestion, learning or work-intensity effects enter other households' values. Such effects must be placed explicitly in the primitives before applying the inequality.

Second, unequal social weights do not permit unweighted transfer cancellation. If household $i$ pays an extra dollar to owner $h$, weighted welfare changes by $\omega_h-\omega_i$ even if the physical allocation is unchanged. More generally, the rent component is
\[
 \sum_{t,j}\beta^t p_{jt}
 \left(\sum_h\omega_h\theta_{hj}q_{jt}
       -\sum_i\omega_ix_{ijt}\right),
\]
where $q_{jt}=\sum_ix_{ijt}$ is occupancy and $\theta_{hj}$ is owner $h$'s share. It vanishes for equal weights with complete ownership coverage. If some owners lie outside the chosen population boundary, their receipts do not cancel inside that boundary; the omitted external-area component must be reported.

A simple diagnostic is a pure rent redistribution: a tenant's annual rent falls by 100, its landlord's income falls by 100, and housing services and other resources do not change. Equal-weight aggregate benefit is zero, even though tenant benefit is positive and property value may fall. Conversely, an annual commuting-resource saving of 100 with unchanged production is a real gain if it is not offset elsewhere. Observing only the same property-price change in two cities cannot distinguish these cases.

\section*{Identification and remaining work}
The optimization is a solution only after $a$, capacities, conversion costs, ownership and terminal values are specified. Baseline observed rents generally do not identify the complete matrix of household-location values or their remote-work counterfactuals. Several value matrices share the same observed optimal assignment and dual rents but assign different values to newly feasible remote configurations. Those observationally equivalent baseline economies can have different $W^*(1)-W^*(0)$.

Empirical resolution requires intervention variation and information about noneligible service workers, selection into remote jobs, productivity, office demand, public finance and conversion. A twenty-year calculation additionally needs a justified transition path and terminal assumption. Taking annual capital gains and then adding the capitalized value of the same future rent stream would count it twice. The present result supplies an exact conditional allocation and accounting benchmark, while leaving those identification and dynamic-equilibrium questions unresolved.

\section{Completion Estimate}
33\%

This is a post hoc, heuristic estimate for the full catalogue target.
The attempt solves a finite urban assignment benchmark over twenty years and proves ownership and rent accounting, conditional conversion optimization, and an option value result. The full remote work welfare target still requires identified productivity and location responses, municipal effects, unequal weighted incidence, conversion implementation, and a credible transition path.

\begin{thebibliography}{9}
\bibitem{review} S. Van Nieuwerburgh (2022), The Remote Work Revolution: Impact on Real Estate Values and the Urban Environment, NBER Working Paper 30662. \url{https://www.nber.org/papers/w30662}.
\bibitem{summary} S. Van Nieuwerburgh (2022), Real Estate Values in the Time of COVID, \emph{NBER Reporter}, number 4. Primary research summary and reference to Working Paper 30662 checked. \url{https://www.nber.org/reporter/2022number4/real-estate-values-time-covid}.
\end{thebibliography}
\end{document}
